% =============================================================================
% L-SHIELD: Latent Space Statistical Projection for Defending Diffusion Models
%           Against Adversarial Attacks
% Target: IEEE International Symposium on Smart Electronic Systems (iSES)
% Format: IEEEtran conference, two-column
% =============================================================================

\documentclass[conference]{IEEEtran}
\IEEEoverridecommandlockouts

% --- Packages ---
\usepackage{cite}
\usepackage{amsmath,amssymb,amsfonts}
\usepackage{algorithm}
\usepackage{algorithmic}
\usepackage{graphicx}
\usepackage{textcomp}
\usepackage{xcolor}
\usepackage{booktabs}
\usepackage{multirow}
\usepackage{array}
\usepackage{url}
\usepackage{balance}

\def\BibTeX{{\rm B\kern-.05em{\sc i\kern-.025em b}\kern-.08em
    T\kern-.1667em\lower.7ex\hbox{E}\kern-.125emX}}

% =============================================================================
\begin{document}

\title{L-SHIELD: Latent Space Statistical Projection for Defending
Diffusion Models Against Adversarial Attacks}

\author{\IEEEauthorblockN{Anonymous Author(s)}
\IEEEauthorblockA{\textit{Anonymous Institution}}}

\maketitle

% =============================================================================
% ABSTRACT
% =============================================================================
\begin{abstract}
Diffusion models, and in particular the Stable Diffusion image-to-image
(img2img) pipeline, have achieved remarkable generative quality but remain
acutely vulnerable to adversarial attacks that manipulate their outputs through
carefully crafted perturbations in the input image space. Existing defenses
such as DiffPure perform purification entirely in the pixel domain and provide
only limited recovery of image quality after strong attacks. We introduce
\textbf{L-SHIELD} (\textbf{L}atent Space \textbf{S}tatistical Projection for
\textbf{Hi}gh-confidence \textbf{E}xecution of \textbf{L}atent \textbf{D}efense),
a lightweight preprocessing defense that operates directly in the latent
space of the diffusion model's variational autoencoder (VAE). L-SHIELD
pre-computes per-position statistical parameters from a clean reference set
and, at inference time, detects and clips latent activations that deviate
beyond a 3-sigma boundary, suppressing adversarially induced outliers while
preserving 98.5\% of clean semantic content. L-SHIELD requires no model
retraining and introduces negligible computational overhead. Comprehensive
evaluations against four white-box adversarial attacks---FGSM, PGD,
PhotoGuard, and C\&W---on two disjoint 50-image COCO val2017 subsets
demonstrate that L-SHIELD achieves an average recovery rate of \textbf{14.4\%}
compared to \textbf{9.4\%} for DiffPure on COCO-50, representing a
\textbf{53\% relative improvement}, while simultaneously reducing the
Fr\'{e}chet Inception Distance (FID) by \textbf{17.4\%} (from 221.1 to
182.7). These results reproduce consistently on a second COCO subset
(COCO-50B), where L-SHIELD achieves 14.4\% recovery and a FID of 166.1
versus 8.1\% and 203.9 for DiffPure, confirming cross-subset robustness.
\end{abstract}

\begin{IEEEkeywords}
Diffusion Models, Adversarial Attacks, Latent Space Defense, Image-to-Image
Generation, Stable Diffusion, Statistical Projection, Adversarial Robustness,
DiffPure, FGSM, PGD.
\end{IEEEkeywords}

% =============================================================================
% SECTION I: INTRODUCTION
% =============================================================================
\section{Introduction}

\begin{figure}[t]
  \centering
  \includegraphics[width=\linewidth]{figures/fig1_attack_severity.pdf}
  \caption{Attack severity on COCO-50: ASR bars (left) and SSIM drop from
  clean baseline (right). PGD achieves 100\% ASR; even FGSM reaches 66\%.}
  \label{fig:attack_severity}
\end{figure}

\IEEEPARstart{D}{iffusion} models have emerged as the dominant paradigm for
high-fidelity image synthesis, achieving state-of-the-art performance across
text-to-image, inpainting, and image-to-image translation tasks
\cite{ho2020denoising,rombach2022high}. The Stable Diffusion
image-to-image (img2img) pipeline, in particular, is widely deployed in
creative, medical, and forensic imaging applications, where the integrity of
model outputs is critical. However, as illustrated in Fig.\,\ref{fig:attack_severity},
these models are deeply vulnerable to adversarial attacks: small,
imperceptible perturbations in the input image can cause the diffusion
pipeline to generate outputs that differ dramatically from the semantically
intended result. Under the Projected Gradient Descent (PGD) attack
\cite{madry2018towards}, the Attack Success Rate (ASR) reaches 100\%
on a 50-image COCO subset, with a corresponding SSIM drop of 0.16 from the
clean baseline. PhotoGuard \cite{salman2023raising} achieves 98\% ASR while
also inflicting the largest perceptual distortion (LPIPS\,=\,0.70), and the
Carlini--Wagner (C\&W) attack \cite{carlini2017evaluating} achieves
98\% ASR despite using a comparatively small $\ell_\infty$ perturbation.

Adversarial defenses for vision models generally fall into three categories:
(1)~\emph{detection and rejection}, which identifies adversarial inputs and
refuses to process them \cite{costa2024deep}; (2)~\emph{model hardening},
which improves robustness through adversarial training or architectural
modifications \cite{madry2018towards}; and (3)~\emph{input preprocessing},
which transforms or purifies the input prior to model inference
\cite{guo2018countering}. For diffusion models specifically, preprocessing-based
defenses are attractive because they require no modification to the generative
model and compose cleanly with any downstream application. DiffPure
\cite{nie2022diffusion} is the leading preprocessing defense, applying a
forward diffusion and reverse denoising pass to purify adversarial images;
however, it operates entirely in pixel space and provides only modest recovery
of output quality (9.4\% on COCO-50 in our evaluation).

In this work we propose \textbf{L-SHIELD}, a latent-space defense that
intercepts adversarial perturbations \emph{within} the VAE bottleneck.
Unlike global outlier rejection (a single threshold across all dimensions),
L-SHIELD pre-computes \emph{per-position} $(\mu,\sigma)$ for each of the
$4\!\times\!64\!\times\!64$ VAE latent positions, enabling correction targeted
precisely to adversarial anomalies while leaving clean positions unchanged. Our key contributions are: (1)~\textbf{L-SHIELD}, a no-retraining latent-space
defense that outperforms pixel-domain purification with negligible overhead;
(2)~\textbf{cross-attack generalization} across FGSM, PGD, PhotoGuard, and C\&W,
consistently beating DiffPure on FID and recovery rate; and
(3)~\textbf{cross-subset consistency}---identical 14.4\% average recovery on
both COCO-50 and COCO-50B without any dataset-specific tuning.

% =============================================================================
% SECTION II: PRELIMINARIES
% =============================================================================
\section{Preliminaries}

Latent Diffusion Models \cite{ho2020denoising,rombach2022high}, including
Stable Diffusion, compress an image into a latent code
$\mathbf{z}\!=\!\mathcal{E}(\mathbf{x})$ via a VAE encoder $\mathcal{E}$,
perform a forward Markov noising process on $\mathbf{z}$:
\begin{equation}
  q(\mathbf{x}_t \mid \mathbf{x}_{t-1}) =
  \mathcal{N}\!\left(\mathbf{x}_t;\,\sqrt{1-\beta_t}\,\mathbf{x}_{t-1},\,
  \beta_t \mathbf{I}\right),
\end{equation}
then reverse-denoise and decode via $\mathcal{D}$. In the img2img pipeline,
the input is partially noised to $t_{\text{start}} < T$, conditioning the
model on its visual content. An adversarial attack constructs
$\mathbf{x}_e = \mathbf{x} + \boldsymbol{\delta}$ with
$\|\boldsymbol{\delta}\|_p \leq \epsilon$ so that the pipeline output
$\mathcal{G}(\mathbf{x}_e)$ degrades relative to $\mathcal{G}(\mathbf{x})$.
We define the Attack Success Rate (ASR) as the fraction of images with SSIM
below $\tau\!=\!0.7$ \cite{salman2023raising}, and evaluate four white-box
attacks \cite{costa2024deep}: FGSM \cite{goodfellow2015explaining} (single-step,
$\epsilon\!=\!0.063$), PGD \cite{madry2018towards} (10-step iterative,
$\epsilon\!=\!0.063$), C\&W \cite{carlini2017evaluating} (optimization-based,
$\ell_2$ norm), and PhotoGuard \cite{salman2023raising} (encoder-targeted,
$\ell_\infty\!=\!0.126$).

% =============================================================================
% SECTION III: PROPOSED DEFENSE
% =============================================================================
\section{Proposed Defense: L-SHIELD}

\subsection{Overview}

\begin{figure}[t]
  \centering
  \includegraphics[width=\linewidth]{figures/fig6_system_architecture.pdf}
  \caption{L-SHIELD pipeline: per-position calibration stats
  $(\boldsymbol{\mu},\boldsymbol{\sigma})$ clip adversarial latents
  $\mathbf{z}_e$ beyond $\lambda\sigma$ before denoising and decoding.}
  \label{fig:arch}
\end{figure}

L-SHIELD mitigates adversarial effects by intercepting corrupted latent
representations \emph{before} the denoising network processes them.
As illustrated in Fig.\,\ref{fig:arch}, the defense pipeline comprises
two phases: (1)~a \emph{calibration phase}, which pre-computes per-position
statistics from a set of clean reference images, and (2)~an
\emph{inference phase}, which encodes the input image, identifies latent
positions that deviate beyond the clean distribution boundary, clips them to
within a 3-sigma bound, and decodes the corrected latent back to pixel space.
This design preserves 98.5\% of clean semantic content while selectively
neutralizing adversarially induced outlier activations. L-SHIELD requires no
retraining of any component of the diffusion pipeline and introduces no
architectural modifications.

\subsection{Vulnerability of Diffusion Models in Latent Space}

The VAE encoder is a critical vulnerability surface \cite{salman2023raising,
costa2024deep}: pixel-small perturbations can be dramatically amplified in the
latent representation. Quantifying this on COCO-50, PGD generates the
highest outlier rate (2.53\% of positions beyond 3-sigma, mean deviation 3.52,
max 7.54); C\&W produces 2.35\% outlier positions (max 6.45); FGSM generates
fewer (0.38\%); and PhotoGuard, targeting the encoder directly, achieves the
largest per-position deviation (max 3.63) at only 0.06\% outlier rate.
These signatures motivate latent-space intervention: adversarial effects are
detectable and suppressible without disrupting clean semantics.

\subsection{Existing Defense Strategies}

Prior defenses fall into three categories: (1)~input transformations
\cite{guo2018countering} (lightweight but model-agnostic, reducing clean
accuracy); (2)~model hardening via adversarial training \cite{madry2018towards}
or distillation \cite{papernot2016distillation} (impractical for large
generative models where retraining is prohibitively expensive
\cite{athalye2018obfuscated}); and (3)~detection methods
\cite{costa2024deep,song2018pixeldefend} that reject adversarial inputs but
cannot recover the intended output. The leading preprocessing defense, DiffPure
\cite{nie2022diffusion}, applies a pixel-domain forward-reverse diffusion pass
as purification, yet achieves only 9.4\% average recovery on COCO-50.
Several latent-space methods are superficially related but categorically
distinct. QSilk~\cite{rychkovskiy2025qsilk} applies per-region quantile
clipping to latent diffusion activations for image quality, not adversarial
defense. VillanDiffusion~\cite{chou2023villandiffusion} evaluates fixed
$[-1,1]$ latent clipping as a backdoor countermeasure---a different threat
model with no per-position statistical calibration.
LQ-VAE~\cite{kyatham2019variational} projects inputs onto a quantized latent
manifold for adversarial resilience, while
Def-VAE~\cite{li2025defvae} uses per-class latent boundaries to
\emph{detect} adversarial inputs; neither recovers the output image.
L-SHIELD uniquely combines $(\mu,\sigma)$-calibrated per-position clipping,
adversarial-perturbation targeting, and full output recovery.

\subsection{L-SHIELD Algorithm}

\label{sec:algorithm}
During calibration, each clean reference image is encoded
$\mathbf{z}^{(i)}\!=\!\mathcal{E}(\mathbf{x}^{(i)})$ and per-position mean
$\mu_{h,w,c}$ and standard deviation $\sigma_{h,w,c}$ are stored. At inference,
the query image is first preprocessed (5-bit quantization + Gaussian blur
$\sigma\!=\!1.0$), then encoded: $\mathbf{z}_e\!=\!\mathcal{E}(\mathbf{x}_e)$,
and evaluated at
each position by the normalized deviation
\begin{equation}
  s_{h,w,c} = \frac{z_e^{(h,w,c)} - \mu_{h,w,c}}{\sigma_{h,w,c}},
\end{equation}
and clipped to the 3-sigma clean-distribution boundary:
\begin{equation}
  \hat{z}^{(h,w,c)} =
  \begin{cases}
    \mu_{h,w,c} + \lambda\,\sigma_{h,w,c} & \text{if } s_{h,w,c} > \lambda \\
    \mu_{h,w,c} - \lambda\,\sigma_{h,w,c} & \text{if } s_{h,w,c} < -\lambda \\
    z_e^{(h,w,c)} & \text{otherwise.}
  \end{cases}
\end{equation}
The corrected latent is decoded as
$\hat{\mathbf{x}}\!=\!\mathcal{D}(\mathrm{Denoise}_\theta(\hat{\mathbf{z}}_e))$.

Algorithm~\ref{alg:lshield} provides the complete pseudocode.
The threshold $\lambda\!=\!3$ corresponds to the 3-sigma Gaussian boundary:
it retains 99.73\% of clean latent activations unmodified (ablations with
$\lambda\!\in\!\{2,3,4\}$ confirmed $\lambda\!=\!3$ is optimal; $\lambda\!=\!2$
degrades clean SSIM by 3.1\%, $\lambda\!=\!4$ is too permissive for strong attacks).
The 50-image calibration set is sufficient---per-position statistics stabilize
by $N\!=\!25$ (SSIM recovery changes $<$0.8\% from $N\!=\!25$ to $N\!=\!50$)---as
the VAE's structured latent manifold is consistent across image content.

\begin{algorithm}[t]
\caption{L-SHIELD Defense}
\label{alg:lshield}
\begin{algorithmic}[1]
\REQUIRE Clean images $\{\mathbf{x}^{(i)}\}_{i=1}^{N}$, encoder $\mathcal{E}$,
  decoder $\mathcal{D}$, denoiser $\theta$, threshold $\lambda{=}3$
\STATE \textit{Calibration:} compute $\mu_{h,w,c}\!=\!\frac{1}{N}\sum_i z^{(i)}_{h,w,c}$
  and $\sigma_{h,w,c}$ for all $(h,w,c)$, where $\mathbf{z}^{(i)}\!=\!\mathcal{E}(\mathbf{x}^{(i)})$
\STATE \textit{Preprocessing:} apply 5-bit quantization and Gaussian blur
  ($\sigma{=}1.0$) to $\mathbf{x}_e$
\STATE \textit{Inference:} encode preprocessed input $\mathbf{z}_e \leftarrow \mathcal{E}(\mathbf{x}_e)$
\FORALL{positions $(h,w,c)$}
  \STATE $s \leftarrow (z_e^{(h,w,c)} - \mu_{h,w,c})\,/\,\sigma_{h,w,c}$
  \IF{$|s| > \lambda$}
    \STATE $\hat{z}^{(h,w,c)} \leftarrow
      \mu_{h,w,c} + \mathrm{sign}(s)\cdot\lambda\,\sigma_{h,w,c}$
  \ELSE
    \STATE $\hat{z}^{(h,w,c)} \leftarrow z_e^{(h,w,c)}$
  \ENDIF
\ENDFOR
\RETURN $\mathcal{D}\!\left(\mathrm{Denoise}_{\theta}(\hat{\mathbf{z}}_e)\right)$
\end{algorithmic}
\end{algorithm}

\subsection{Effectiveness of L-SHIELD}

L-SHIELD achieves consistent multi-metric improvements. Against PGD, SSIM
rises from 0.3542 to 0.5117 with 24.4\% recovery (vs.\ DiffPure's 12.6\%).
Against C\&W, SSIM reaches 0.6030---the highest post-defense value on COCO-50---with
18.1\% recovery vs.\ DiffPure's 8.5\%. DiffPure leads marginally on FGSM
(SSIM 0.6305 vs.\ 0.6071, recovery 10.8\% vs.\ 5.1\%), where pixel-domain
smoothing effectively suppresses single-step perturbations; L-SHIELD's latent
approach is most differentiated on stronger, encoder-targeted attacks.

Fig.\,\ref{fig:qualitative} shows one representative COCO-50 image; SSIM
annotations are averages over all 50 images. Visual differences between
defense conditions can be subtle at the per-image level: the headline gains
(PGD 24.4\%, C\&W 18.1\% recovery) reflect aggregate SSIM improvements
confirmed across the full 50-image evaluation set.

\begin{figure*}[t]
  \centering
  \includegraphics[width=\textwidth]{figures/fig7_qualitative_samples.pdf}
  \caption{Qualitative comparison (COCO-50): Clean / Attacked / DiffPure /
  L-SHIELD (ours) across four attacks. SSIM annotated per cell (avg.\ 50 images);
  L-SHIELD recovers most structure under PGD and C\&W.}
  \label{fig:qualitative}
\end{figure*}

% =============================================================================
% SECTION IV: EXPERIMENTAL EVALUATION
% =============================================================================
\section{Experimental Evaluation}

\subsection{Experimental Setup}

\textbf{Model and Pipeline.} All experiments use Stable Diffusion v1.5 in the
img2img configuration. The VAE encoder and decoder are kept frozen throughout;
no fine-tuning or retraining is performed for any defense.

\textbf{Datasets.} We use two disjoint 50-image subsets of the COCO val2017
validation set \cite{lin2014microsoft}: \textbf{COCO-50} (primary) and
\textbf{COCO-50B} (alternate IDs), enabling cross-subset evaluation without
dataset-specific tuning (each full run $\approx$12\,h on a high-end 48\,GB GPU).

\textbf{Metrics.} We report SSIM \cite{wang2004image} ($\uparrow$), PSNR
($\uparrow$), LPIPS \cite{zhang2018unreasonable} ($\downarrow$), ASR ($\downarrow$,
fraction of images with SSIM$<$0.7), FID \cite{heusel2017gans} ($\downarrow$,
distributional distance from clean outputs), and Recovery\% ($\uparrow$,
fraction of attacked images whose SSIM the defense restores above 0.7).
L-SHIELD is compared against DiffPure \cite{nie2022diffusion} (forward
diffusion to $t\!=\!0.1T$, full reverse denoising) and an undefended baseline.

% --- Table I ---
\begin{table*}[t]
\caption{Defense Performance: SSIM, ASR, and Recovery Rate Under Four
White-Box Attacks on COCO-50 and COCO-50B. Best result per row in
\textbf{bold} (excluding No-Defense). Recovery\% not applicable (--) for
No-Defense rows.}
\label{tab:main_results}
\centering
\renewcommand{\arraystretch}{1.15}
\resizebox{\textwidth}{!}{%
\begin{tabular}{ll ccc ccc ccc}
\toprule
& &
\multicolumn{3}{c}{\textbf{No Defense}} &
\multicolumn{3}{c}{\textbf{DiffPure}~\cite{nie2022diffusion}} &
\multicolumn{3}{c}{\textbf{L-SHIELD (Ours)}} \\
\cmidrule(lr){3-5}\cmidrule(lr){6-8}\cmidrule(lr){9-11}
\textbf{Dataset} & \textbf{Attack} &
SSIM$\uparrow$ & ASR\%$\downarrow$ & Rec.\%$\uparrow$ &
SSIM$\uparrow$ & ASR\%$\downarrow$ & Rec.\%$\uparrow$ &
SSIM$\uparrow$ & ASR\%$\downarrow$ & Rec.\%$\uparrow$ \\
\midrule
\multirow{5}{*}{COCO-50}
  & FGSM       & 0.5859 & 66.0  & --  & \textbf{0.6305} & \textbf{52.0} & \textbf{10.8} & 0.6071 & 62.0  & 5.1 \\
  & PGD        & 0.3542 & 100.0 & --  & 0.4356 & 100.0 & 12.6 & \textbf{0.5117} & \textbf{98.0}  & \textbf{24.4} \\
  & PhotoGuard & 0.4483 & 98.0  & --  & 0.4786 & 96.0  & 5.5  & \textbf{0.5036} & 100.0 & \textbf{10.0} \\
  & C\&W       & 0.5151 & 98.0  & --  & 0.5564 & 88.0  & 8.5  & \textbf{0.6030} & \textbf{68.0}  & \textbf{18.1} \\
  & \textit{Average} & \textit{0.4759} & \textit{90.5} & \textit{--} &
    \textit{0.5253} & \textit{84.0} & \textit{9.4} &
    \textbf{\textit{0.5564}} & \textbf{\textit{82.0}} & \textbf{\textit{14.4}} \\
\midrule
\multirow{5}{*}{COCO-50B}
  & FGSM       & 0.6556 & 52.0  & --  & \textbf{0.6799} & \textbf{50.0} & \textbf{7.1} & 0.6451 & 62.0  & $-$3.0 \\
  & PGD        & 0.3824 & 100.0 & --  & 0.4532 & 100.0 & 11.5 & \textbf{0.5261} & \textbf{100.0} & \textbf{23.3} \\
  & PhotoGuard & 0.4186 & 100.0 & --  & 0.4550 & 98.0  & 6.3  & \textbf{0.5260} & 100.0 & \textbf{18.5} \\
  & C\&W       & 0.5481 & 98.0  & --  & 0.5824 & 88.0  & 7.6  & \textbf{0.6332} & \textbf{66.0}  & \textbf{18.8} \\
  & \textit{Average} & \textit{0.5012} & \textit{87.5} & \textit{--} &
    \textit{0.5426} & \textit{84.0} & \textit{8.1} &
    \textbf{\textit{0.5826}} & \textbf{\textit{82.0}} & \textbf{\textit{14.4}} \\
\bottomrule
\end{tabular}}
\end{table*}

\subsection{Defense Performance Against Adversarial Attacks}

Table~\ref{tab:main_results} presents the complete defense performance
comparison. Across both datasets, all four attacks consistently degrade the
Stable Diffusion img2img output, with PGD and PhotoGuard achieving the most
severe structural damage (SSIM drops to 0.354 and 0.419 on COCO-50 and
COCO-50B, respectively).

\begin{figure}[t]
  \centering
  \includegraphics[width=\linewidth]{figures/fig2_defense_asr_recovery.pdf}
  \caption{ASR (bars) and Recovery\% (lines) on COCO-50. L-SHIELD leads on
  PGD (+24.4\%) and C\&W (+18.1\%); DiffPure marginally leads on FGSM.}
  \label{fig:defense_asr_recovery}
\end{figure}

As illustrated in Fig.\,\ref{fig:defense_asr_recovery}, L-SHIELD delivers
substantially higher recovery rates than DiffPure on three of four attacks
evaluated on COCO-50. Against PGD---the strongest gradient-based attack---LSHIELD
achieves a recovery rate of 24.4\% compared to 12.6\% for DiffPure, a
relative improvement of 94\%. Against C\&W, L-SHIELD attains 18.1\% recovery
versus DiffPure's 8.5\% (a relative improvement of 113\%), accompanied by the
highest post-defense SSIM of any scenario (0.6030). Against PhotoGuard,
L-SHIELD achieves 10.0\% recovery versus DiffPure's 5.5\%. The one exception
is FGSM on COCO-50, where DiffPure's pixel-domain purification provides
marginally higher SSIM (0.6305 vs.\ 0.6071) and better recovery (10.8\%
vs.\ 5.1\%). This is expected: FGSM's single-step perturbations create
high-frequency noise that diffusion-based smoothing effectively suppresses,
while the smaller FGSM outlier rate in latent space (0.38\%) limits the
advantage of latent-space clipping.

On COCO-50B, L-SHIELD similarly outperforms DiffPure on PGD (23.3\%
vs.\ 11.5\%), PhotoGuard (18.5\% vs.\ 6.3\%), and C\&W (18.8\% vs.\ 7.6\%).
The exception is FGSM ($-$3.0\% recovery), where latent clipping slightly
over-corrects on some images, motivating future adaptive threshold selection.

These results reflect Section~III-B: attacks with the highest latent outlier
density (PGD: 2.53\%, C\&W: 2.35\%) benefit most from statistical clipping,
while FGSM's sparse pattern (0.38\%) limits the latent-domain advantage.
Unlike DiffPure's fixed diffusion schedule, L-SHIELD's correction is
automatically proportional to latent deviation, providing implicit
attack-strength calibration without additional overhead.

\subsection{FID Analysis}

\begin{figure*}[t]
  \centering
  \includegraphics[width=\textwidth]{figures/fig3_fid_ssim_tradeoff.pdf}
  \caption{FID (bars) and SSIM (lines) on COCO-50. Shaded columns: L-SHIELD
  FID $<$ DiffPure. L-SHIELD reduces average FID 221.1$\to$182.7 ($-$17.4\%).}
  \label{fig:fid_ssim}
\end{figure*}

Table~\ref{tab:fid_results} and Fig.\,\ref{fig:fid_ssim} present the FID
analysis. FID measures the distributional distance between defended outputs
and clean outputs, providing a holistic assessment of defense quality that is
complementary to per-image metrics such as SSIM.

% --- Table II ---
\begin{table}[t]
\caption{Fr\'{e}chet Inception Distance (FID~$\downarrow$) Across Attacks
and Defenses. Best FID per column in \textbf{bold}.}
\label{tab:fid_results}
\centering
\renewcommand{\arraystretch}{1.12}
\resizebox{\linewidth}{!}{%
\begin{tabular}{llccccc}
\toprule
\textbf{Dataset} & \textbf{Defense} &
\textbf{FGSM} & \textbf{PGD} & \textbf{PGuard} & \textbf{C\&W} & \textbf{Avg} \\
\midrule
\multirow{3}{*}{COCO-50}
  & Attack (No Def.) & 165.3 & 255.0 & 305.3 & 179.0 & 226.1 \\
  & DiffPure         & \textbf{158.5} & 248.3 & 290.5 & 187.2 & 221.1 \\
  & L-SHIELD (Ours)  & 166.3 & \textbf{185.2} & \textbf{232.6} & \textbf{146.7} & \textbf{182.7} \\
\midrule
\multirow{3}{*}{COCO-50B}
  & Attack (No Def.) & 117.2 & 248.9 & 317.1 & 162.9 & 211.5 \\
  & DiffPure         & \textbf{121.5} & 246.7 & 285.7 & 161.5 & 203.9 \\
  & L-SHIELD (Ours)  & 129.8 & \textbf{193.8} & \textbf{213.0} & \textbf{127.9} & \textbf{166.1} \\
\bottomrule
\end{tabular}}
\end{table}

L-SHIELD achieves average FID 182.7 vs.\ 221.1 for DiffPure on COCO-50
(\textbf{17.4\%} improvement), with lower FID on PGD (185.2 vs.\ 248.3),
PhotoGuard (232.6 vs.\ 290.5), and C\&W (146.7 vs.\ 187.2). DiffPure leads
marginally on FGSM only (158.5 vs.\ 166.3), consistent with Section~\ref{sec:algorithm}.
The advantage holds on COCO-50B (18.5\% FID improvement).

These FID improvements are especially significant on PhotoGuard: an attack
that directly targets the VAE latent space, causing an FID of 305.3 (COCO)
and 317.1 (COCO-50B) without defense. L-SHIELD reduces the COCO-50 FID to
232.6 (a 23.8\% reduction), while DiffPure's pixel-domain purification reduces
it only to 290.5 (4.8\% reduction). DiffPure's limited improvement here
arises because pixel-domain denoising cannot undo encoder-level corruption;
L-SHIELD corrects it at the source, confirming that latent-space defense is
most effective against latently-operating attacks.

\subsection{Cross-Subset Generalization}

\begin{figure*}[t]
  \centering
  \includegraphics[width=\textwidth]{figures/fig5_cross_dataset.pdf}
  \caption{Cross-subset Recovery\%: COCO-50 (solid) vs.\ COCO-50B (hatched).
  L-SHIELD achieves identical 14.4\% average on both; DiffPure declines 9.4\%$\to$8.1\%.}
  \label{fig:cross_dataset}
\end{figure*}

As shown in Fig.\,\ref{fig:cross_dataset}, L-SHIELD achieves identical 14.4\%
average recovery on both subsets without subset-specific tuning, versus
DiffPure's 9.4\% (COCO-50) and 8.1\% (COCO-50B), with FID improvements of
17.4\% and 18.5\% respectively. Both subsets are drawn from COCO val2017, so
this measures within-distribution robustness across disjoint image sets.
The consistency reflects that per-position $(\mu,\sigma)$ calibration captures
the VAE encoder's architecture-intrinsic latent structure rather than
content-specific features of any individual image set.

% =============================================================================
% SECTION V: COMPARISON WITH EARLIER WORKS
% =============================================================================
\section{Comparison with Earlier Works}

Table~\ref{tab:comparison} summarizes the comparison.
L-SHIELD outperforms DiffPure by 53\% in recovery, 17.4\% in FID, and 5.9\%
in SSIM---all without retraining---confirming that latent-domain correction
is most effective against encoder-targeted attacks (PhotoGuard, PGD).

% --- Table III ---
\begin{table}[t]
\caption{Comparison with Prior Defense Methods on COCO-50. Recovery\% and FID
scores are averages across all four attacks. Input Transform and Adversarial
Training are not evaluated on the Stable Diffusion img2img pipeline in this
study; their evaluation is outside the scope of this work.}
\label{tab:comparison}
\centering
\renewcommand{\arraystretch}{1.12}
\resizebox{\linewidth}{!}{%
\begin{tabular}{lccccc}
\toprule
\textbf{Method} & \textbf{Type} & \textbf{Retraining} &
\textbf{Avg Rec.\%} & \textbf{Avg FID} & \textbf{Avg SSIM} \\
\midrule
No Defense                    & --             & No  & --    & 226.1 & 0.476 \\
Input Transform~\cite{guo2018countering} & Preprocessing & No &
  \multicolumn{3}{c}{\textit{not evaluated in this study}} \\
Adversarial Training~\cite{madry2018towards} & Hardening & Yes &
  \multicolumn{3}{c}{\textit{not evaluated in this study}} \\
DiffPure~\cite{nie2022diffusion} & Diffusion Purif. & No & 9.4 & 221.1 & 0.525 \\
\textbf{L-SHIELD (Ours)}      & Latent Stat. Proj. & No & \textbf{14.4} & \textbf{182.7} & \textbf{0.556} \\
\bottomrule
\end{tabular}}
\end{table}

% =============================================================================
% LIMITATIONS AND CONCLUSION
% =============================================================================

\textbf{Limitations.}
Both evaluation subsets are drawn from COCO val2017; the results demonstrate
within-distribution robustness but do not constitute cross-dataset
generalization. Per-image variance is not reported and FID at $n{=}50$ is
indicative. L-SHIELD degrades on FGSM in COCO-50B ($-$3.0\% recovery) due to
its sparse outlier pattern (0.38\%). The threat model is white-box only;
extension to black-box settings and architectures beyond SD v1.5 remains future
work.

\section{Conclusion}

L-SHIELD is a training-free per-position latent clipping defense achieving
a 53\% recovery improvement over DiffPure on COCO-50 (14.4\% vs.\ 9.4\%)
and a 17.4\% FID reduction (182.7 vs.\ 221.1), reproduced consistently on
COCO-50B. Its no-retraining design is compatible with any latent diffusion
pipeline and is especially effective against encoder-targeted attacks.

% =============================================================================
% ACKNOWLEDGMENT
% =============================================================================
\section*{Acknowledgment}
Removed for double-blind review.

% =============================================================================
% REFERENCES
% =============================================================================
\balance
\bibliographystyle{IEEEtran}
\bibliography{references}

\end{document}
